\section{Routing (Exercise 3.2)}\label{sec:routing}
\noindent Responsible: Person B. Notebook: \code{notebooks/ex32\_routing.ipynb}.

\subsection{Road network from OpenStreetMap (2 points)}

\paragraph{Method.}
\todo{Describe the OpenStreetMap extract \cite{geofabrik-osm} (region file, date), how it was cropped to the study area
(e.g.\ with osmium-tool and a buffered boundary), how the road network was extracted and converted into a routable
graph (pyrosm \cite{pyrosm} or OSMnx \cite{osmnx}; network type, edge attributes used as routing weight, e.g.\ length
or travel time, handling of one-way streets and of the largest connected component), and the number of nodes and edges
(\cref{tab:network-stats}).}

\begin{table}[H]
  \centering
  \caption{\todo{Caption: size of the routing network.}}
  \label{tab:network-stats}
  \begin{tabular}{@{}lr@{}}
    \toprule
    Quantity & Value \\
    \midrule
    Nodes & \todo{value} \\
    Edges & \todo{value} \\
    Total length [km] & \todo{value} \\
    \bottomrule
  \end{tabular}
\end{table}

\begin{figure}[H]
  \centering
  \figplaceholder{network_map}
  \caption{\todo{Caption: road network of the study area.}}
  \label{fig:network-map}
\end{figure}

\paragraph{Result.}
\todo{Describe \cref{fig:network-map}.}

\subsection{Routing of the trips and flow per link (2 points)}

\paragraph{Method.}
\todo{Describe how origin and destination coordinates are mapped to the nearest network node (CRS, distance of the snap),
which shortest-path algorithm and weight are used, how the number of trips traversing each edge is counted, how trips
without a path are handled, and the runtime.}

\begin{figure}[H]
  \centering
  \figplaceholder{edge_flows_map}
  \caption{\todo{Caption: number of trips per network link.}}
  \label{fig:edge-flows-map}
\end{figure}

\paragraph{Result.}
\todo{Report the number of routed trips, trips without a route, and the most loaded links, and comment on
\cref{fig:edge-flows-map} (major roads, bottlenecks, bridges, effect of the border of the study area).}

\paragraph{Problems encountered.}
\todo{Document problems, e.g.\ the extract is larger than the study area, unreachable nodes, long runtimes,
memory, the network ends at the border of the department although trips cross it.}
